\documentclass[a4paper,12pt]{article}
\usepackage[T2A]{fontenc}
\usepackage[utf8]{inputenc}
\usepackage[russian]{babel}
\usepackage{amsmath,amssymb,mathtools}
\renewcommand{\sfdefault}{cmss}
\renewcommand{\familydefault}{\sfdefault}
\usepackage{icomma}
\usepackage[margin=22mm,top=22mm,bottom=23mm]{geometry}
\usepackage{microtype}
\usepackage{tikz}
\usetikzlibrary{arrows.meta,positioning,shapes.geometric,calc}
\usepackage{pgfplots}
\pgfplotsset{compat=1.18}
\usepackage{booktabs,tabularx}
\usepackage{enumitem}
\usepackage{fancyhdr}
\usepackage{parskip}
\usepackage{xcolor}
\usepackage{setspace}
\usepackage[hidelinks]{hyperref}
\definecolor{accent}{HTML}{087B80}
\definecolor{darkaccent}{HTML}{07545A}
\definecolor{textgray}{HTML}{353B40}
\definecolor{muted}{HTML}{66737A}
\definecolor{lightaccent}{HTML}{E5F2F0}
\definecolor{softline}{HTML}{B4D5D4}
\color{textgray}
\setstretch{1.12}
\setlength{\parindent}{0pt}
\setlength{\parskip}{6pt}
\setlist[itemize]{leftmargin=1.5em,itemsep=2pt,topsep=3pt}
\setlist[enumerate]{leftmargin=1.6em,itemsep=6pt,topsep=3pt}
\pagestyle{fancy}
\fancyhf{}
\fancyhead[L]{\small\color{muted}Математика \quad\textbullet\quad 25.10.2026}
\fancyhead[R]{\small\color{muted}Матвей Горбачев}
\fancyfoot[C]{\small\color{muted}\thepage}
\renewcommand{\headrulewidth}{0.25pt}
\renewcommand{\headrule}{\hbox to\headwidth{\color{softline}\leaders\hrule height \headrulewidth\hfill}}
\setlength{\headheight}{26pt}
\newcommand{\sectionline}[1]{\vspace{3mm}{\Large\bfseries\color{darkaccent}#1}\par\vspace{1mm}{\color{softline}\hrule height 0.6pt}\vspace{2mm}}
\newcommand{\subline}[1]{{\large\bfseries\color{darkaccent}#1}\par}
\newcommand{\key}[1]{{\color{darkaccent}\bfseries #1}}
\newcommand{\task}[2]{\item \textbf{#1} #2}
\pgfplotsset{lessonplot/.style={width=\linewidth,height=5.15cm,axis lines=middle,axis line style={-Stealth,semithick},xlabel={$x$},ylabel={$y$},xlabel style={at={(ticklabel* cs:1)},anchor=north west},ylabel style={at={(ticklabel* cs:1)},anchor=south east},grid=both,major grid style={draw=gray!20},minor tick num=0,ticklabel style={font=\small},label style={font=\small},every axis plot/.append style={very thick,accent},samples=180,clip=false}}
\begin{document}
\thispagestyle{empty}
\begin{center}
\vspace*{29mm}
{\color{accent}\bfseries\Large МАТЕМАТИКА \quad / \quad ПОДГОТОВКА К ЕГЭ}\par
\vspace{22mm}
{\color{darkaccent}\fontsize{31}{38}\selectfont\bfseries Чек-лист}\par
\vspace{11mm}
{\fontsize{21}{28}\selectfont\bfseries Производная и графики}\par
\vspace{3mm}
{\fontsize{16}{22}\selectfont Знак производной, монотонность, экстремумы}\par
\vspace{17mm}
{\color{accent}\rule{42mm}{1pt}}\par
\vspace{15mm}
{\large Для Матвея Горбачева}\par
\vspace{3mm}
{\normalsize 25 октября 2026 года}\par
\vfill
{\color{textgray}\normalsize Лёвин Артём Александрович}\par
{\color{muted}\small эксклюзивно для Матвея Горбачева}\par
\vspace{12mm}
\begin{tikzpicture}[x=1cm,y=1cm]
\draw[accent!55,line width=1.3pt] (-7.1,0) .. controls (-4.9,1.1) and (-2.8,-0.9) .. (0,0.0) .. controls (2.6,0.9) and (4.9,-0.3) .. (7.1,0.75);
\draw[softline,line width=0.6pt] (-7.1,-0.18) .. controls (-4.0,0.7) and (-1.7,-1.0) .. (1.2,-0.3) .. controls (3.7,0.2) and (5.2,-0.5) .. (7.1,0.3);
\end{tikzpicture}
\end{center}
\clearpage
\sectionline{1. Сначала определите, чей график дан}
\begin{tabularx}{\linewidth}{@{}p{0.27\linewidth}X X@{}}
\toprule
& \textbf{Дан график $y=f(x)$} & \textbf{Дан график $y=f'(x)$} \\
\midrule
\textbf{Что смотрим?} & Куда идёт кривая слева направо & Где кривая лежит относительно оси $Ox$ \\
\textbf{Возрастание $f$} & Кривая поднимается & $f'(x)>0$ (над осью) \\
\textbf{Убывание $f$} & Кривая опускается & $f'(x)<0$ (под осью) \\
\textbf{Экстремум $f$} & Переход роста к убыванию или наоборот & Смена знака $f'$ \\
\bottomrule
\end{tabularx}
\vspace{2mm}
\subline{Три главные связи}
\[
\begin{aligned}
&f'(x)>0 &&\Longrightarrow\quad f\text{ возрастает};\\
&f'(x)<0 &&\Longrightarrow\quad f\text{ убывает};\\
&f'(x_0)=0 &&\Longrightarrow\quad x_0\text{ --- стационарная точка}.
\end{aligned}
\]
Последние две строки особенно важны при чтении графика производной: \key{сам график $f'$ может убывать, пока функция $f$ возрастает}, если $f'$ остаётся положительной.

\subline{Как определить максимум и минимум}
\begin{center}
\begin{tikzpicture}[font=\small,>=Stealth]
\node[align=center] at (-4.1,0.65) {\bfseries максимум};
\node[align=center] at (3.9,0.65) {\bfseries минимум};
\node at (-5.2,-0.15) {$+$}; \node at (-3.0,-0.15) {$-$};
\draw[->,thick,accent] (-4.85,-0.7)--(-3.35,-0.7);\node at (-4.1,-0.7) [above] {$f'$};
\node at (-4.1,-1.27) {$f:\ \nearrow\ \searrow$};
\node at (2.8,-0.15) {$-$}; \node at (5.1,-0.15) {$+$};
\draw[->,thick,accent] (3.15,-0.7)--(4.65,-0.7);\node at (3.9,-0.7) [above] {$f'$};
\node at (3.9,-1.27) {$f:\ \searrow\ \nearrow$};
\end{tikzpicture}
\end{center}
\textbf{Важно.} У дифференцируемой функции во внутренней точке экстремума производная равна нулю, но \emph{не всякий нуль производной является экстремумом}. Например, у $f(x)=x^3$ в точке $x=0$ имеем $f'(0)=0$, хотя функция продолжает возрастать. В точке излома производная может не существовать.

\sectionline{2. Четыре вопроса к любой задаче}
\begin{enumerate}[label=\textcolor{accent}{\bfseries\arabic*.}]
\item \textbf{Чей график дан?} Найдите в условии $f(x)$ или $f'(x)$.
\item \textbf{Каково условие отбора?} Например, $f'(x)<0$, точка максимума, наименьшее значение.
\item \textbf{Что это значит для данного графика?} Переведите условие на язык роста, убывания или знака производной.
\item \textbf{Что записать?} Число точек, сумму целых $x$, длину промежутка, абсциссу или значение функции.
\end{enumerate}
\clearpage
\sectionline{3. Типовые вопросы ЕГЭ}
\subline{Целые абсциссы}
Если требуется количество целых $x$ при $f'(x)<0$, отмечайте \emph{целые абсциссы} на нужных промежутках. Координата $y$ может быть любой. Точки, где $f'(x)=0$, в строгое неравенство не входят. Выколотые точки и точки вне области определения также исключаются.

\textbf{Пример.} Пусть график производной расположен ниже оси на интервале $(-3;2)$. Тогда
\[
-3<x<2,\quad x\in\mathbb{Z}\quad\Longrightarrow\quad x=-2,-1,0,1.
\]
\textbf{Количество:} $4$. \quad \textbf{Сумма:} $-2$. \quad \textbf{Длина интервала:} $2-(-3)=5$.

\subline{Экстремумы по графику $f'$}
\begin{itemize}
\item Знак $f'$ меняется с $+$ на $-$: у $f$ \key{точка максимума}.
\item Знак $f'$ меняется с $-$ на $+$: у $f$ \key{точка минимума}.
\item Знак $f'$ не меняется: из одного лишь $f'(x_0)=0$ экстремум не следует.
\end{itemize}

\subline{Наибольшее и наименьшее значения на отрезке}
Если $f$ непрерывна на $[a;b]$ и $f'(x)>0$ на $(a;b)$, функция возрастает на $[a;b]$: минимум достигается при $x=a$, максимум --- при $x=b$. Если $f'(x)<0$, всё наоборот.

\textbf{Пример из типа задач занятия.} На отрезке $[3;8]$ производная положительна. Следовательно, $f$ возрастает, и \textbf{наименьшее значение} достигается при $x=3$. Если спрашивают \emph{в какой точке}, ответ --- $3$, а если \emph{чему равно наименьшее значение}, нужны данные о самой функции.

Если производная меняет знак внутри отрезка, для \emph{наибольшего/наименьшего значения} может потребоваться сравнить значения $f$ на концах и в критических точках. По одному знаку $f'$ нельзя всегда определить, какое из нескольких локальных минимумов является наименьшим.

\sectionline{4. Что легко перепутать}
\begin{itemize}
\item Нуль $f'$ означает \emph{горизонтальную касательную} при существовании производной, но не обязательно экстремум.
\item \textbf{Точка максимума} --- значение $x_0$; \textbf{максимум функции} --- значение $f(x_0)$.
\item При работе с графиком $f'$ смотрим \emph{высоту относительно $Ox$}, а не то, растёт ли сама кривая $f'$.
\item При подсчёте целых учитываем строгие/нестрогие границы и выколотые точки.
\item Сумма целых абсцисс, их количество и длина промежутка --- три разных ответа.
\end{itemize}
\clearpage
\sectionline{Алгоритм: как читать график и получать ответ}
\vspace{7mm}
\begin{center}
\begin{tikzpicture}[
  font=\small,>=Stealth,
  start/.style={ellipse,draw=darkaccent,line width=.75pt,fill=lightaccent,minimum width=4.3cm,minimum height=0.85cm,align=center},
  choice/.style={diamond,aspect=2.3,draw=darkaccent,line width=.8pt,fill=lightaccent,inner sep=4pt,align=center,text width=2.7cm},
  flowstep/.style={rounded corners=3pt,draw=softline,line width=.85pt,fill=white,text width=4.0cm,minimum height=1.08cm,align=center},
  main/.style={rounded corners=3pt,draw=softline,line width=.85pt,fill=white,text width=8.3cm,minimum height=0.88cm,align=center},
  link/.style={->,line width=.8pt,darkaccent}
]
\node[start] (start) at (0,0) {Прочитайте условие задачи};
\node[choice] (q) at (0,-2.0) {Чей график дан?};
\node[flowstep] (left) at (-3.2,-4.3) {График $f$:\\определите участки роста и убывания, характер экстремумов};
\node[flowstep] (right) at (3.2,-4.3) {График $f'$:\\определите интервалы $f'>0$, $f'<0$ и смены знака};
\node[main] (filter) at (0,-7.0) {Переведите вопрос на язык графика:\\какие точки или промежутки подходят?};
\node[main] (bounds) at (0,-8.55) {Уточните границы отрезка, нули, выколотые точки};
\node[main] (answer) at (0,-10.10) {Выберите нужное действие: количество, сумма, длина, $x$ или $f(x)$};
\node[start] (finish) at (0,-11.65) {Запишите ответ};
\draw[link] (start)--(q);
\draw[link] (q.west)-| node[pos=.28,above] {$f$} (left.north);
\draw[link] (q.east)-| node[pos=.28,above] {$f'$} (right.north);
\draw[link] (left.south)|-(filter.west);
\draw[link] (right.south)|-(filter.east);
\draw[link] (filter)--(bounds);
\draw[link] (bounds)--(answer);
\draw[link] (answer)--(finish);
\end{tikzpicture}
\end{center}
\clearpage
\sectionline{Тренировка. Путь А --- 10 заданий}
\textit{Решайте по рисункам. Все графики непрерывны и дифференцируемы на показанных промежутках; масштаб по оси $x$ указан на рисунках. Ответы приведены отдельно.}

\subline{График A: $y=f(x)$ на $[-2;2]$}
\begin{center}
\begin{tikzpicture}
\begin{axis}[lessonplot,width=.74\linewidth,height=5.5cm,xmin=-2.3,xmax=2.3,ymin=-3.5,ymax=3.5,xtick={-2,-1,0,1,2},ytick={-2,0,2},ylabel={$f(x)$}]
\addplot[domain=-2:2] {x^3-3*x};
\addplot[only marks,mark=*,mark size=1.4pt,accent] coordinates {(-2,-2) (2,2)};
\end{axis}
\end{tikzpicture}
\end{center}
\begin{enumerate}[label=\textcolor{accent}{\bfseries\arabic*.},itemsep=5pt]
\item Сколько \emph{целых} $x\in[-2;2]$ удовлетворяют условию $f'(x)<0$?
\item Сколько точек на интервале $(-2;2)$ удовлетворяют условию $f'(x)=0$?
\item Укажите абсциссу точки \emph{максимума} функции внутри интервала $(-2;2)$.
\end{enumerate}
\vspace{1mm}
\textit{Подсказка к самостоятельной проверке: обратите внимание, что здесь изображена именно $f(x)$.}
\clearpage
\sectionline{Тренировка. Продолжение}
\subline{График B: $y=f'(x)$ на $[-4;5]$}
\begin{center}
\begin{tikzpicture}
\begin{axis}[lessonplot,width=.86\linewidth,height=6.4cm,xmin=-4.4,xmax=5.4,ymin=-6.6,ymax=5.8,xtick={-4,-3,-2,-1,0,1,2,3,4,5},ytick={-6,-3,0,3},ylabel={$f'(x)$}]
\addplot[domain=-4:5] {(x+3)*(x-1)*(x-4)/8};
\addplot[only marks,mark=*,mark size=1.5pt,accent] coordinates {(-4,-5) (5,4)};
\end{axis}
\end{tikzpicture}
\end{center}
\begin{enumerate}[label=\textcolor{accent}{\bfseries\arabic*.},start=4,itemsep=7pt]
\item Сколько \emph{внутренних точек минимума} имеет функция $f$ на $(-4;5)$?
\item Укажите абсциссу \emph{внутренней точки максимума} функции $f$.
\item Найдите \emph{сумму всех целых} $x\in[-4;5]$, при которых $f'(x)<0$.
\item Найдите \emph{длину наибольшего промежутка убывания} $f$ в пределах $[-4;5]$.
\end{enumerate}
\clearpage
\sectionline{Тренировка. Завершение}
\subline{График C: $y=f'(x)$ на $[-4;4]$}
\begin{center}
\begin{tikzpicture}
\begin{axis}[lessonplot,width=.84\linewidth,height=6.1cm,xmin=-4.4,xmax=4.4,ymin=-5,ymax=13.6,xtick={-4,-3,-2,-1,0,1,2,3,4},ytick={-4,0,4,8,12},ylabel={$f'(x)$}]
\addplot[domain=-4:4] {x^2-4};
\addplot[only marks,mark=*,mark size=1.5pt,accent] coordinates {(-4,12) (4,12)};
\end{axis}
\end{tikzpicture}
\end{center}
\begin{enumerate}[label=\textcolor{accent}{\bfseries\arabic*.},start=8,itemsep=10pt]
\item Сколько целых $x\in[-4;4]$ удовлетворяют условию $f'(x)<0$?
\item В какой точке отрезка $[-1;1]$ функция $f$ принимает \emph{наименьшее значение}?
\item Укажите абсциссу \emph{внутренней точки минимума} $f$ на $(-4;4)$.
\end{enumerate}
\vspace{8mm}
\textbf{Самопроверка перед ответами.} Для каждого задания назовите: чей график изображён; какое условие требуется; как оно выглядит на графике; какую величину записывают в ответ.
\clearpage
\sectionline{Ответы к тренировочному блоку}
\vspace{6mm}
\renewcommand{\arraystretch}{1.45}
\begin{tabularx}{\linewidth}{@{}c c X@{}}
\toprule
\textbf{№} & \textbf{Ответ} & \textbf{Короткая проверка} \\
\midrule
1 & $1$ & Убывание на $(-1;1)$; целая абсцисса: $0$. \\
2 & $2$ & Горизонтальные касательные при $x=-1$ и $x=1$. \\
3 & $-1$ & В точке $x=-1$ рост сменяется убыванием. \\
4 & $2$ & $f'$ меняет знак $-\to+$ при $x=-3$ и $x=4$. \\
5 & $1$ & $f'$ меняет знак $+\to-$ при $x=1$. \\
6 & $1$ & Подходят целые $x=-4,\,2,\,3$; сумма $1$. \\
7 & $3$ & Длиннейший интервал убывания --- от $1$ до $4$. \\
8 & $3$ & Подходят $x=-1,\,0,\,1$. \\
9 & $1$ & На $[-1;1]$ имеем $f'<0$, поэтому минимум справа. \\
10 & $2$ & При $x=2$ производная меняет знак $-\to+$. \\
\bottomrule
\end{tabularx}
\vspace{14mm}
{\color{darkaccent}\bfseries Главный вывод занятия.} В задачах по графикам производной решающее значение имеет \emph{знак} $f'$; в задачах по графику функции --- направление изменения $f$. Проверяйте формулировку ответа: количество, сумма, длина, абсцисса и значение функции не взаимозаменяемы.
\end{document}
